\documentclass[a4paper,11pt]{article}
\usepackage[utf8]{inputenc}
\usepackage[T1]{fontenc}
\usepackage[ngerman]{babel}
\usepackage{amsmath,amssymb}
\usepackage{xcolor}
\usepackage{colortbl}
\usepackage{array}
\usepackage{tikz}
\usepackage{enumitem}
\usepackage[margin=2cm]{geometry}
\definecolor{bgdark}{HTML}{1E1E1E}
\definecolor{fgtext}{HTML}{E6E6E6}
\definecolor{gridcol}{HTML}{4A4A4A}
\definecolor{accent}{HTML}{6CB6FF}
\definecolor{loescol}{HTML}{7FD18C}
\pagecolor{bgdark}
\color{fgtext}
\arrayrulecolor{fgtext}
\setlength{\parindent}{0pt}
\newcommand{\karo}[1]{\par\noindent\begin{tikzpicture}\draw[step=5mm,gridcol,very thin](0,0) grid (\linewidth,#1);\end{tikzpicture}\par\vspace{2mm}}
\newcommand{\aufgabe}[2]{\vspace{4mm}\par\noindent{\large\color{accent}\textbf{Aufgabe #1}}\hfill{\color{gridcol}\normalsize #2}\par\vspace{1mm}}
\newcommand{\titel}[1]{\begin{center}{\LARGE\color{accent}\textbf{Optimierungsverfahren}}\\[3pt]{\large #1}\end{center}\vspace{1mm}\noindent\textbf{Name:}\ \underline{\hspace{6cm}}\hfill\textbf{Datum:}\ \underline{\hspace{4cm}}\\[3pt]{\color{gridcol}\rule{\linewidth}{0.4pt}}\vspace{2mm}}

\begin{document}
\titel{Übungsaufgaben --- Blatt 05: Mehrzieloptimierung \& Methodenübersicht}

\aufgabe{1 --- Dominanz \& Pareto-Menge}{mittel}
Gegeben sind sechs Punkte im Zielraum (beide Ziele $f_1,f_2$ werden \textbf{minimiert}):
\[
P_1=(2,6),\ P_2=(3,4),\ P_3=(5,3),\ P_4=(4,5),\ P_5=(6,2),\ P_6=(4,7).
\]
\begin{enumerate}[label=(\alph*)]
  \item Geben Sie für jeden Punkt an, welche anderen Punkte er dominiert.
  \item Welche Punkte sind \textbf{Pareto-optimal}?
\end{enumerate}
\karo{6cm}

\aufgabe{2 --- Non-Dominated-Sorting-Rang}{mittel}
Gegeben sind die Koordinaten von sechs Punkten im Zielraum mit \textbf{drei} Zielfunktionen
(alle werden \textbf{minimiert}):
\[
P_1=(1,2,3),\ P_2=(2,1,2),\ P_3=(3,2,1),\ P_4=(2,2,2),\ P_5=(3,3,3),\ P_6=(1,3,1).
\]
Bestimmen Sie für jeden Punkt den Non-Dominated-Sorting-Rang.
\karo{6.5cm}

\aufgabe{3 --- Crowding Distance}{mittel}
Im Zielraum (2 Ziele, \textbf{Minimierung}) liegen die nicht-dominierten Punkte
\[
A=(1,7),\ B=(2,5),\ C=(3,4),\ D=(5,3),\ E=(7,1).
\]
\begin{enumerate}[label=(\alph*)]
  \item Berechnen Sie die Crowding Distance jedes Punktes.
  \item Welche \textbf{vier} Punkte würde man anhand der Crowding Distance auswählen?
\end{enumerate}
\karo{7cm}

\aufgabe{4 --- Hypervolumen}{mittel--schwer}
Gegeben die (nicht-dominierte) Pareto-Front $\{(1,7),(2,5),(3,4),(5,3),(7,1)\}$ (Minimierung beider Ziele)
und der Referenzpunkt $\vec r=(8,8)$. Berechnen Sie das \textbf{Hypervolumen} (die von der Front gegenüber
$\vec r$ dominierte Fläche).
\karo{7cm}

\aufgabe{5 --- NSGA-II Selektion}{mittel}
Eine Population enthält fünf Individuen mit Rang $1$ und zwei mit Rang $2$. Die Crowding Distances der
Rang-1-Individuen sind:
\[
a:\infty,\quad b:0{,}4,\quad c:1{,}0,\quad d:\infty,\quad e:0{,}7.
\]
Es sollen $4$ Individuen für die nächste Generation überleben. Welche werden ausgewählt und nach welcher
Regel? Begründen Sie.
\karo{5cm}

\aufgabe{6 --- Eigenschaften der Verfahren}{mittel}
Ordnen Sie die folgenden Verfahren nach den Eigenschaften ein: \emph{global?}, \emph{gradientenbasiert?},
\emph{deterministisch?}, \emph{Anzahl Ziele (1/mehrere)}.
\[
\text{Gradientenverfahren},\ \text{Simplex (LP)},\ \text{KKT},\ \text{Nelder-Mead},\
\text{DIRECT},\ \text{(1+1)-ES},\ \text{Simulated Annealing},\ \text{NSGA-II}.
\]
\karo{9cm}

\end{document}
